% Archived BK-CMA material. This file is intentionally not input by main.tex.
% It preserves the Schnorr and lattice key-recovery results for possible later use.
% It assumes the notation and generic tunable-lossy-function definition of main.tex.

\subsection*{Archived preliminaries}

For a signature scheme $\SigScheme=(\SigKeys,\SigSign,\SigVf)$, the key-recovery game gives the adversary the verification key and an adaptive signing oracle, and the adversary wins by returning the signing key:
\begin{center}
\fbox{
\begin{minipage}{2.8in}
\ExperimentHeader{Game $\gameBKCMA{\SigScheme}$}

\begin{oracle}{$\Initialize$}
\item $(\vkey,\skey) \getsr \SigKeys$
\item Return $\vkey$
\end{oracle}
\ExptSepSpace
\begin{oracleC}{$\SignOracle(m)$}{$m \in \MsgSp$}
\item $\sigma\getsr\SigSign(\skey,m)$
\item Return $\sigma$
\end{oracleC}
\ExptSepSpace
\begin{oracle}{$\Finalize(\skey')$}
\item Return $(\skey'=\skey)$
\end{oracle}
\end{minipage}
}
\end{center}
The BK-CMA advantage of $\advA$ is
$$
\AdvBKCMA{\SigScheme}{\advA}
:=
\Pr\big[\gameBKCMA{\SigScheme}(\advA)\Rightarrow1\big].
$$
For Schnorr signatures we identify $\skey$ with the secret exponent $x$.

We say that the injective mode of a tunable lossy function is \emph{universal} if, for every $r,r'\in[M]$, the marginal distributions of $K$ in
$$
(K,\mathsl{td})\getsr\TLFKg(\usecp,r,\mathsf{inj})
\qquad\text{and}\qquad
(K,\mathsl{td})\getsr\TLFKg(\usecp,r',\mathsf{inj})
$$
are identical.

A \emph{strongly efficiently enumerable ELF} is the special case with a universal injective mode in which, for every polynomial time bound $t$ and inverse polynomial $\epsilon$, there is a polynomial $q$ such that every adversary of time at most $t(\secp)$ has $\AdvTLF{\TLF,r}{\advA}\le\epsilon(\secp)$ for every $r\ge q(\secp)$. For polynomial $r$, the functions $\nu_{\mathsf{inj}}(\secp,r)$ and $\nu_{\mathsf{enum}}(\secp,r)$ are negligible, and $\tau(\secp,r)$ is polynomial in $\secp$ and $r$.

For an ELF with domain size $M$, the injective generator runs Zhandry's key generator $\ELFKg$ on $(M,M)$ and the lossy generator with bound $r$ runs it on $(M,r)$. Zhandry constructs such an ELF from exponential DDH, or more generally from a constant-$k$ exponential $k$-linear assumption~\cite{C:Zhandry16}. The construction first builds a bounded-adversary lossy function and then iterates it at several auxiliary security levels. The enumeration and security properties above follow from Theorem~3.9, Lemma~3.6, and Corollary~3.15 of the full version~\cite{EPRINT:Zhandry16}.

\section{Key-Recovery Security of Schnorr Signatures} \label{sec-lf-kr}

We instantiate the hash function with a tunable lossy function in injective mode, followed by a collision-resistant outer hash into $\Zp$. The outer hash supplies the scalar challenge required by Schnorr and, when the injective-mode failure probability is negligible, makes the deployed composite collision resistant. Its collision resistance is not used in the key-recovery reduction. The main theorem is concrete and holds for every image bound $r$. A strongly efficiently enumerable ELF gives the asymptotic corollary: the reduction runs in polynomial time, so ordinary DL hardness suffices in the Schnorr group.

\subsection{The Construction} \label{sec-lf-constr}

Let $\GG$ be a cyclic group of prime order $p$ with generator $g$. Let $\TLF$ be a tunable lossy function with algorithms $\TLFKg,\TLFEv,\TLFEnum$. Its domain is $\bits^{\ellin}$ and its range is $\Ygrave$. Recall from Section~\ref{sec-lf} that $M=2^{\ellin}$, and fix an injective encoding $\enc{\cdot}\colon\GG\times\bits^{\msglen}\to\bits^{\ellin}$.

Let $\CRHF=(\CRHFKg,\CRHFEv)$ be a keyed family from $\Ygrave$ to $\Zp$. Algorithm $\CRHFKg$ on input $\usecp$ outputs a key $k$, and $\CRHFEv$ is deterministic and efficient. We require collision resistance: for every PPT adversary $\advA$,
$$
\Pr\left[
\begin{array}{l}
k\getsr\CRHFKg(\usecp);\ (y_0,y_1)\getsr\advA(k):\\
y_0\ne y_1\ \land\ \CRHFEv(k,y_0)=\CRHFEv(k,y_1)
\end{array}
\right]
$$
is negligible. A collision-resistant hash with $\mu=\Omega(\secp)$ output bits, followed by the canonical embedding $\bits^\mu\hookrightarrow\Zp$, gives such a family when $2^\mu\le p$.

\begin{construction} \label{constr-lf}
Fix an image parameter $r\in[M]$. The keyed hash family $\HF^{\mathsf{tlf}}_r=(\HFKg,\HFEv)$ is defined as follows. Algorithm $\HFKg$ on input $X\in\GG$ computes
$$
(K,\mathsl{td})\getsr\TLFKg(\usecp,r,\mathsf{inj})
\qquad\text{and}\qquad
k\getsr\CRHFKg(\usecp),
$$
and returns $\hk=(K,k)$. Algorithm $\HFEv$ on inputs $\hk=(K,k)$ and $(R,m)\in\GG\times\bits^{\msglen}$ returns
$$
\CRHFEv\bigl(k,\TLFEv(K,\enc{R,m})\bigr).
$$
\end{construction}

\paragraph*{Discussion.} Key generation ignores $X$, and lossy keys are used only in the proof. The outer hash maps the possibly structured range $\Ygrave$ to the scalar required by Schnorr. When the injective-mode failure probability is negligible, its collision resistance makes the composite hash collision resistant by Proposition~\ref{prop-lf-cr}. This property is separate from key recovery: the simulator needs only to apply the outer map to the enumerated lossy image. Without the condition, the allowed outer family could be constant. If its value is $c$, anyone can forge on any message by choosing $s\getsr\Zp$ and returning $(g^sX^{-c},s)$. We impose collision resistance to exclude such degenerate hash functions and to make the deployed composite collision resistant, not to justify the BK-CMA reduction.

\begin{proposition} \label{prop-lf-cr}
If $\CRHF$ is collision resistant and $\nu_{\mathsf{inj}}(\secp,r)$ is negligible, then $\HF^{\mathsf{tlf}}_r$ is collision resistant.
\end{proposition}

\begin{proof}
Condition on $\TLFEv(K,\cdot)$ being injective. Suppose an adversary returns distinct inputs $(R_0,m_0)$ and $(R_1,m_1)$ with equal hash values. Since the encoding and $\TLFEv(K,\cdot)$ are injective,
$$
y_0:=\TLFEv(K,\enc{R_0,m_0})
\qquad\text{and}\qquad
y_1:=\TLFEv(K,\enc{R_1,m_1})
$$
are distinct. Their outer hashes are equal, giving a collision for $\CRHF$. The excluded event has probability at most $\nu_{\mathsf{inj}}(\secp,r)$.
\end{proof}

\subsection{Main Result} \label{sec-lf-thm}

\begin{theorem} \label{thm-bk}
Let $\advA$ be an adversary against BK-CMA of $\SchScheme[\GG,\HF^{\mathsf{tlf}}_r]$ that runs in time $t_{\advA}$ and makes at most $q_s$ signing queries. For every $r\in[M]$, there are a distinguisher $\advD$ and a DL algorithm $\advB$ such that
$$
\AdvBKCMA{\SchScheme[\GG,\HF^{\mathsf{tlf}}_r]}{\advA}
\le
\AdvTLF{\TLF,r}{\advD}
+\AdvDL{\GG,g}{\advB}
+\nu_{\mathsf{enum}}(\secp,r)
+q_s e^{-\secp}.
$$
Adversary $\advD$ runs in time $t_{\advA}+q_s\poly(\secp)$, and $\advB$ runs in time
$$
t_{\advA}+\tau(\secp,r)+(q_s+1)\secp r\poly(\secp).
$$
\end{theorem}

The proof replaces the signing oracle by rejection sampling over the enumerated lossy image. Conditioned on acceptance, the simulated signature has exactly the honest distribution.

The games are in Figure~\ref{fig-bk-games}. In the figure, write
$$
H(R,m)=\CRHFEv(k,\TLFEv(K,\enc{R,m})).
$$
Game $\Gm_0$ is the BK-CMA game. Game $\Gm_1$ changes the tunable lossy function to lossy mode at the same value $r$. Game $\Gm_2$ enumerates the lossy image once and replaces honest signing by rejection sampling; if $C$ is empty, which occurs only within the enumeration failure event, the oracle returns $\bot$.

\begin{figure}[t]
\twoCols{0.43}{0.52}{
\ExperimentHeader{Game $\mathsf{G}[\mathsf{mode},b]$}

\begin{oracle}{$\Initialize$}
\item $x\getsr\Zp\;;\;X\gets g^x$
\item $k\getsr\CRHFKg(\usecp)$
\item $(K,\mathsl{td})\getsr\TLFKg(\usecp,r,\mathsf{mode})$
\item If $b=1$:
\item \quad $S\getsr\TLFEnum(\mathsl{td})$
\item \quad $C\gets\{\CRHFEv(k,y):y\in S\}$
\item Return $(X,K,k)$
\end{oracle}
\ExptSepSpace
\begin{oracle}{$\Finalize(x')$}
\item Return $(x'=x)$
\end{oracle}
}
{
\begin{oracle}{$\SignOracle(m)$}
\item If $b=0$:
\item \quad $r'\getsr\Zp\;;\;R\gets g^{r'}$
\item \quad $c\gets H(R,m)\;;\;s\gets(r'+cx)\bmod p$
\item \quad Return $(R,s)$
\item For $i=1,\ldots,\secp|C|$:
\item \quad $s\getsr\Zp\;;\;c\getsr C$
\item \quad $R\gets g^sX^{-c}$
\item \quad If $H(R,m)=c$: Return $(R,s)$
\item Return $\bot$
\end{oracle}
}
\vspace{-12pt}
\caption{Games for the proof of Theorem~\ref{thm-bk}. We set $\Gm_0=\mathsf{G}[\mathsf{inj},0]$, $\Gm_1=\mathsf{G}[\mathsf{lossy},0]$, and $\Gm_2=\mathsf{G}[\mathsf{lossy},1]$.}\label{fig-bk-games}
\vspace{-3pt}\hrulefill
\vspace{-1ex}
\end{figure}

We first prove that the rejection-sampling oracle is exact.

\begin{lemma} \label{lem-sim}
Fix $x\in\Zp$, $X=g^x$, a function $H \colon \GG \times \bits^{\msglen} \to \Zp$, a message $m$, and a finite set $C$ with $\Image(H) \subseteq C$. For $(c,s)\in C\times\Zp$, call $(c,s)$ good if
$$
H\bigl(g^sX^{-c},m\bigr)=c.
$$
Then $\varphi(c,s):=g^sX^{-c}$ is a bijection from the set of good pairs to $\GG$.
\end{lemma}

\begin{proof}
For $R\in\GG$, define
$$
c_R:=H(R,m)
\qquad\text{and}\qquad
s_R:=\mathsf{dlog}_g(R)+c_Rx\pmod p.
$$
The hypothesis $\Image(H) \subseteq C$ gives $c_R\in C$. We also have
$$
g^{s_R}X^{-c_R}=R,
$$
so $(c_R,s_R)$ is good and $\varphi(c_R,s_R)=R$.

Conversely, let $(c,s)$ be good and set $R=\varphi(c,s)$. Goodness gives $c=c_R$. The equation $R=g^sX^{-c}$ then gives $s=s_R$. Thus the good pairs are exactly $\{(c_R,s_R):R\in\GG\}$, and $\varphi$ maps each such pair to $R$.
\end{proof}

A trial in game $\Gm_2$ is uniform on $C\times\Zp$. Apply Lemma~\ref{lem-sim} with $H(R,m) := \CRHFEv(k,\TLFEv(K,\enc{R,m}))$ and the set $C$ of the game; on executions where enumeration succeeds, $\Image(H) \subseteq C$ holds. The lemma gives exactly $p$ good pairs, so a trial accepts with probability $1/|C|$. Conditioned on acceptance, $R$ is uniform in $\GG$, and the returned response is $s=\mathsf{dlog}_g(R)+c_Rx\bmod p$. This is exactly an honest Schnorr signature. The probability that all $\secp|C|$ trials fail is at most $e^{-\secp}$.

\begin{proof}[of Theorem~\ref{thm-bk}]
The distinguisher $\advD$, on input $K$, samples $x\getsr\Zp$ and $k\getsr\CRHFKg(\usecp)$. It runs $\advA$ on $(g^x,K,k)$, answers signing queries honestly, and returns 1 exactly when $\advA$ outputs $x$. The two input distributions for $\advD$ give games $\Gm_0$ and $\Gm_1$. Hence
$$
|\Pr[\Gm_0\Rightarrow1]-\Pr[\Gm_1\Rightarrow1]|
\le\AdvTLF{\TLF,r}{\advD}.
$$

Next condition on successful enumeration in $\Gm_2$. By Lemma~\ref{lem-sim}, each signing query has exactly the distribution in $\Gm_1$ unless its rejection loop fails. Therefore
$$
|\Pr[\Gm_1\Rightarrow1]-\Pr[\Gm_2\Rightarrow1]|
\le\nu_{\mathsf{enum}}(\secp,r)+q_se^{-\secp}.
$$

Finally, algorithm $\advB$, on input a DL challenge $X$, samples
$$
(K,\mathsl{td})\getsr\TLFKg(\usecp,r,\mathsf{lossy})
\qquad\text{and}\qquad
k\getsr\CRHFKg(\usecp),
$$
and runs $\advA$ on $(X,K,k)$. It answers signing queries as in $\Gm_2$ and returns the value output by $\advA$. This perfectly simulates $\Gm_2$, so
$$
\Pr[\Gm_2\Rightarrow1]=\AdvDL{\GG,g}{\advB}.
$$
The image is enumerated once. Constructing $C$ takes $r\poly(\secp)$ time, and each signing query uses at most $\secp r$ trials. This gives the claimed running times and completes the proof.
\end{proof}

\paragraph*{Why the construction uses two modes.} The second mode is necessary because the rejection-sampling simulator is also an attack. The oracle of $\Gm_2$ uses only the verification key, and by Lemma~\ref{lem-sim} each of its trials produces a valid signature with probability exactly $1/|C|$. An adversary that runs the same loop on a message of its choice therefore obtains a signature on it. This gives the following bound, which holds for every hash function and uses nothing about how it is built. In the statement, a plain function $H$ is taken as a keyed hash family whose key generation outputs the empty key.

\begin{proposition} \label{prop-lossy-forge}
Let $H \colon \GG \times \bits^{\msglen} \to \Zp$ be efficient, and suppose a set $C \supseteq \Image(H)$ of size at most $r$ can be computed in time $\tau$. Then there is an adversary that runs in time $\tau + \secp r \poly(\secp)$, makes no signing queries, and wins the EUF-CMA game against $\SchScheme[\GG,H]$ with probability at least $1 - e^{-\secp}$.
\end{proposition}

\begin{proof}
The adversary fixes any message $m^*$ and repeats the loop of the $\SignOracle$ of Figure~\ref{fig-bk-games} on $m^*$, which uses only $X$ and $H$. By Lemma~\ref{lem-sim} each trial succeeds with probability $1/|C|$, and by the argument preceding the proof of Theorem~\ref{thm-bk} the probability that all $\secp|C|$ trials fail is at most $e^{-\secp}$. Any output of the loop is a valid signature on $m^*$, and no signing query was made.
\end{proof}

This has two consequences. First, a deployed hash function with a small image gives a scheme that is universally forgeable in time about $r$. The attack cost is governed by the image bound and not by the size of the group, so a large group does not compensate for a small image. With the quasi-polynomial $r = 2^{(\log \secp)^2}$ at $\secp = 128$, for instance, a forger runs in about $2^{49}$ steps whatever group is used. The deployed scheme therefore uses the injective mode, whose image is large and admits no such attack. The lossy mode occurs only inside the proof, and mode indistinguishability transfers the conclusion of the lossy game to the deployed scheme.

Second, and independently of our construction, Proposition~\ref{prop-lossy-forge} shows that lossiness alone cannot give unforgeability. Any image small enough to rejection-sample against is small enough to forge against, in the same time. So the technique of this section stops at key recovery for a structural reason, and unforgeability requires a different tool. That tool is the subject of Section~\ref{sec-io-uf}. The conversion function there has small preimages, so its image is large, and the loop above is too slow to be a simulation or an attack.

For a strongly efficiently enumerable ELF, in the sense stated after Definition~\ref{def-lf}, the injective mode does not depend on the image parameter, so the distribution of $\HF^{\mathsf{tlf}}_r$ is the same for every $r$. We write $\HF^{\mathsf{tlf}}$ for it.

\begin{corollary} \label{cor-bk}
Suppose $\TLF$ is a strongly efficiently enumerable ELF and DL is hard against PPT algorithms in $\GG$. Then $\SchScheme[\GG,\HF^{\mathsf{tlf}}]$ is BK-CMA secure.
\end{corollary}

\begin{proof}
Suppose a PPT adversary $\advA$ has nonnegligible BK-CMA advantage. Fix an inverse polynomial $\epsilon$ such that this advantage is at least $4\epsilon$ for infinitely many security parameters. The distinguisher $\advD$ of Theorem~\ref{thm-bk} is PPT, and its code is independent of $r$: it receives $K$ and never uses the image bound. Let $q$ be the polynomial given by the ELF guarantee for $\advD$ and $\epsilon$, and set $r(\secp)=\lceil q(\secp)\rceil$. Then $\AdvTLF{\TLF,r}{\advD}\le\epsilon$. The injective-mode distribution does not depend on this choice of $r$, so the left side of the theorem is still the advantage against the honest scheme.

For this fixed polynomial $r$, algorithm $\advB$ is PPT. Its DL advantage is negligible. The enumeration-failure and truncation terms are also negligible. For all sufficiently large parameters the right side is less than $4\epsilon$, a contradiction.
\end{proof}

\paragraph*{Instantiation and comparison.} Zhandry constructs a strongly efficiently enumerable ELF from exponential DDH, or constant-$k$ exponential $k$-linear~\cite{C:Zhandry16}. Its honest key contains auxiliary components at several security levels. The honest distribution is fixed. For a given attack, the lossy hybrid selects a component whose security level is large enough. Corollary~\ref{cor-bk} therefore uses ordinary DL in the Schnorr group; neither the DL reduction nor the Schnorr group must have subexponential security. The auxiliary ELF construction assumes exponential DDH, and the proof chooses $r$ as a function of the attack.

The proof of Theorem~\ref{thm-bk} is pointwise in $r$, so it also applies to injective and lossy distributions defined at one fixed value of $r$. The one-level DDH construction of Peikert and Waters~\cite[Sections~5.2--5.3]{STOC:PeiWat08} gives an illustrative example. Let its auxiliary group have order $q'\le M$, set $r=q'$, and retain the vector that defines the lossy mode as the enumeration trapdoor. The lossy image is then enumerable in time $q'\poly(\secp)$. Ignoring constants and lower-order terms, the generic-group estimates have the form
$$
\poly(\ellin)\frac{t_{\advD}^2}{q'}
+
\frac{\bigl(t_{\advA}+q_s\secp q'\poly(\secp)\bigr)^2}{p}.
$$
This estimate illustrates the tradeoff: increasing $q'$ also makes enumeration more expensive. We do not claim the displayed expression as a full concrete-security proposition. The ELF cascade instead uses ordinary DL in the Schnorr group and exponential DDH in its auxiliary groups.

Paillier and Vergnaud~\cite{AC:PaiVer05} prove BK-CMA security for Schnorr with any hash function under OMDL. Their result covers \emph{any} hash function, while ours covers a particular keyed family. Our group assumption is ordinary DL. The hash construction additionally uses exponential DDH for the ELF and collision resistance for the separate hash property of Proposition~\ref{prop-lf-cr}; the BK-CMA reduction does not use the latter. OMDL is interactive; all of our assumptions are non-interactive.

\section{Key-Recovery Security of Lyubashevsky Signatures}

The lattice result uses the following search version of LWE. The advantage of $\advA$ against search $\LWE_{q,n,m',t,B}$ is
$$
\AdvSLWE{q,n,m',t,B}{\advA}
:=
\Pr\big[\mat{A}\mat{V}_1' + \mat{V}_2' = \mat{T} \,\land\, \mat{V}_1' \in [\pm B]^{m'\times t} \,\land\, \mat{V}_2' \in [\pm B]^{n\times t}\big] \;,
$$
where $\mat{A} \getsr \Zq^{n\times m'}$, $\mat{V}_1 \getsr [\pm B]^{m'\times t}$, $\mat{V}_2 \getsr [\pm B]^{n\times t}$, $\mat{T} \gets \mat{A}\mat{V}_1+\mat{V}_2$, and $(\mat{V}_1',\mat{V}_2') \getsr \advA(\mat{A},\mat{T})$. The winner returns any consistent short pair, not necessarily the planted one.

\subsection{The Construction} \label{sec-dil-bk}

The hash family composes the tunable lossy function with an outer hash into $\Zq^n$ and then applies $\Fq_q$ to reach the challenge space. Let $\TLF$ be a tunable lossy function with domain $\bits^{\ellin}$, where $\ellin \ge nk + \msglen$ so that pairs $(\mat{w},m)$ embed by an injective encoding $\enc{\cdot}'' \colon \Zq^n \times \bits^{\msglen} \to \bits^{\ellin}$, and let $(\CRHFKg,\CRHFEv)$ be a collision-resistant keyed family from $\Ygrave$ to $\Zq^n$, as in Section~\ref{sec-lf-constr} but with range $\Zq^n$.

\begin{construction} \label{constr-dil-tlf}
Fix an image parameter $r\in[M]$. The keyed hash family $\HF^{\mathsf{ltlf}}_r=(\HFKg,\HFEv)$ for $\Zq^n$ and message space $\bits^{\msglen}$ is defined as follows. Algorithm $\HFKg$ on input $(\mat{A},\mat{T})$ computes $(K,\mathsl{td})\getsr\TLFKg(\usecp,r,\mathsf{inj})$ and $k'\getsr\CRHFKg(\usecp)$ and returns $\hk=(K,k')$; we write $k'$ because $k$ is the modulus exponent in this part. Algorithm $\HFEv$ on inputs $\hk$ and $(\mat{w},m)$ returns
$$
\Fq_q\!\left(\CRHFEv\bigl(k',\TLFEv(K,\enc{\mat{w},m}'')\bigr)\right) \;.
$$
\end{construction}

\paragraph*{Discussion.} Key generation ignores $(\mat{A},\mat{T})$, and the real scheme uses only injective TLF keys. The lossy mode appears only in the proof, where its enumerable image lets the reduction rejection-sample signatures without the short secret. The outer hash maps the possibly structured TLF range into $\Zq^n$. As in Section~\ref{sec-lf-kr}, collision resistance excludes a constant deployed hash but is not used in the key-recovery reduction.

\subsection{Main Result} \label{sec-dil-bk-thm}

The following lemma is the lattice analogue of Lemma~\ref{lem-sim}, and it is exact. It contains the reason rejection sampling and lossy simulation compose: the secret is recoverable from an accepted response, so accepted attempts correspond one-to-one with the pairs the simulator samples. We write $\mat{y} = (\mat{y}_1, \mat{y}_2)$, $\mat{z} = (\mat{z}_1, \mat{z}_2)$, and $\mat{S}c := (\mat{S}_1c, \mat{S}_2c)$, so that $\mat{z} = \mat{y} + \mat{S}c$ componentwise.

\begin{lemma} \label{lem-dil-sim}
Fix $\mat{A}$, short $\mat{S}_1,\mat{S}_2$ with entries in $[\pm1]$, $\mat{T}=\mat{A}\mat{S}_1+\mat{S}_2$, a function $H \colon \Zq^n \times \bits^{\msglen} \to \chal_q$, a message $m$, and a finite set $C \subseteq \bits^{\ell}$ with $\Image(H) \subseteq C$. Call a pair $(c,\mat{z})$ with $\mat{z}=(\mat{z}_1,\mat{z}_2)$ \emph{good} if
$$
\|\mat{z}_1\|_\infty \le \Bclip,
\qquad
\|\mat{z}_2\|_\infty \le \Bclip,
\qquad
H\big(\mat{A}\mat{z}_1+\mat{z}_2-\mat{T}c,\ m\big)=c \;.
$$
Consider the honest attempt: $\mat{y} \getsr [\pm B]^{\colsA+n}$, $c \gets H(\mat{A}\mat{y}_1+\mat{y}_2, m)$, $\mat{z}_i \gets \mat{y}_i+\mat{S}_ic$, accepted if both norms are at most $\Bclip$. Consider the simulated attempt: $c \getsr C$, $\mat{z}$ uniform on the clipped boxes, accepted if $(c,\mat{z})$ is good. Then the map $\mat{y} \mapsto (c,\mat{z})$ is a bijection from accepted honest attempts to good pairs, and both attempts, conditioned on acceptance, output a uniformly distributed good pair.
\end{lemma}

\begin{proof}
The honest attempt lands in the good set: acceptance bounds the norms, and $\mat{A}\mat{z}_1+\mat{z}_2-\mat{T}c = \mat{A}\mat{y}_1+\mat{y}_2 + (\mat{A}\mat{S}_1+\mat{S}_2)c - \mat{T}c = \mat{A}\mat{y}_1+\mat{y}_2$, so the good-pair condition is the definition of $c$. The map is injective because $\mat{y} = \mat{z}-\mat{S}c$ recovers the attempt. It is surjective: given a good pair, set $\mat{y} \gets \mat{z}-\mat{S}c$; then $\|\mat{y}_i\|_\infty \le \Bclip + \ell \le B$, so $\mat{y}$ is in the attempt space, the commitment $\mat{A}\mat{y}_1+\mat{y}_2$ equals $\mat{A}\mat{z}_1+\mat{z}_2-\mat{T}c$ by the same computation, the good-pair condition makes the attempt's challenge equal $c$, and the attempt is accepted with response $\mat{z}$. Each accepted attempt arises from exactly one uniformly chosen $\mat{y}$, so the honest conditional distribution is uniform on good pairs. The simulated attempt is uniform on $C \times$ the clipped boxes, and its acceptance event is exactly the good set, because every good pair has $c \in \Image(H) \subseteq C$; conditioning gives the uniform distribution on the same set.
\end{proof}

\paragraph*{Discussion.} Two counting facts make the simulation efficient, and both are independent of $H$. First, every $\mat{y}$ with $\|\mat{y}_i\|_\infty \le \Bclip-\ell$ is accepted regardless of the challenge, since the shift $\mat{S}_ic$ moves it by at most $\ell$; hence the honest attempt accepts with probability at least $\big(\tfrac{2(\Bclip-\ell)+1}{2B+1}\big)^{\colsA+n} \ge 1/2$ under $B \ge 4\ell(\colsA+n)$, and the number of good pairs is at least half the number of clipped responses. Second, therefore, the simulated attempt accepts with probability at least $1/(2|C|)$, so $2\secp|C|$ trials fail with probability at most $e^{-\secp}$.

\begin{theorem} \label{thm-dil-bk}
Let $\advA$ be an adversary against BK-CMA of $\DSScheme[\HF^{\mathsf{ltlf}}_r]$ that runs in time $t_{\advA}$ and makes at most $q_s$ signing queries. For every $r\in[M]$, there are a distinguisher $\advD$ and an algorithm $\advB$ such that
$$
\AdvBKCMA{\DSScheme[\HF^{\mathsf{ltlf}}_r]}{\advA}
\le
\AdvTLF{\TLF,r}{\advD}
+\AdvSLWE{q,n,\colsA,\ell,1}{\advB}
+\nu_{\mathsf{enum}}(\secp,r)
+q_s\big(e^{-\secp}+2^{-\secp}\big) \;.
$$
Adversary $\advD$ runs in time $t_{\advA}+q_s\poly(\secp)$, and $\advB$ runs in time $t_{\advA}+\tau(\secp,r)+(q_s+1)\secp\, r\poly(\secp)$.
\end{theorem}

\begin{proof}
The games mirror Theorem~\ref{thm-bk}. Game $\Gm_0$ is BK-CMA. Game $\Gm_1$ generates the TLF key in lossy mode at the same $r$; the distinguisher $\advD$ samples $\mat{A},\mat{S}_1,\mat{S}_2$ and the outer-hash key itself, signs honestly, and outputs 1 iff $\advA$ returns the signing key, so the switch costs $\AdvTLF{\TLF,r}{\advD}$. Game $\Gm_2$ computes $S\getsr\TLFEnum(\mathsl{td})$ and $C\gets\{\Fq_q(\CRHFEv(k',y)):y\in S\}$ once, and answers each signing query by up to $2\secp|C|$ simulated attempts of Lemma~\ref{lem-dil-sim}, returning $\bot$ if all fail. In lossy mode every hash value lies in $C$, except with probability $\nu_{\mathsf{enum}}(\secp,r)$, so the lemma applies: conditioned on producing an output, each query is answered with exactly the distribution of $\Gm_1$. The output distributions differ only through the two failure events, honest signing exhausting $\secp$ attempts and simulation exhausting $2\secp|C|$ trials, which the counting facts above bound by $2^{-\secp}$ and $e^{-\secp}$ per query. Finally, algorithm $\advB$ receives a search-LWE challenge $(\mat{A},\mat{T})$ with $\ell$ columns at box $1$, which has exactly the distribution of the public key, samples the lossy hash key itself, answers signing queries as in $\Gm_2$ without any secret, and forwards the matrix pair from $\advA$'s output key. If $\advA$ wins BK-CMA, that pair is the planted $(\mat{S}_1,\mat{S}_2)$, which is a valid search-LWE solution. The reduction in fact needs less. It succeeds whenever $\advA$ outputs any pair $(\mat{S}_1',\mat{S}_2')$ of matrices with entries in $[\pm1]$ satisfying $\mat{A}\mat{S}_1'+\mat{S}_2' = \mat{T}$, whether or not it equals the planted one. Theorem~\ref{thm-dil-bk} therefore also holds for the stronger notion in which the adversary wins by producing any such pair, which is the more natural notion here, since the key is not determined by the public key. Hence $\Pr[\Gm_2\Rightarrow1]\le\AdvSLWE{q,n,\colsA,\ell,1}{\advB}$, and the running-time claims follow by inspection.
\end{proof}

\paragraph*{Discussion.} The interpretation matches Section~\ref{sec-lf-kr}: with this hash family, the signing oracle gives no help toward key recovery beyond the public key itself, at the price of assuming search LWE hard for time $r \cdot \mathrm{poly}$. The theorem is generic in the tunable lossy function. A lattice-based instantiation would make every ingredient lattice-based. The candidate is the LWE-based construction of Peikert and Waters at a fixed quasi-polynomial image bound, with a rounded output; Section~\ref{sec-conclusion} sketches the candidate argument for enumerability and a universal injective mode, and records what remains to verify. If that argument is formalized, the theorem will assume search LWE against quasi-polynomial time. No obfuscation occurs anywhere in this section, and the reductions are straight-line, without rewinding, so we expect the result to hold against quantum adversaries that make classical signing queries. \adam{Ojaswi/Carter: Two items: (i) verify the counting-fact constants and the $2\secp|C|$ trial bound against the write-up of Theorem~\ref{thm-bk}; (ii) the BK-CMA win condition of Figure~\ref{fig-eufcma} is exact key equality, which is a special case of the search-LWE solution condition, so the last reduction step is sound as stated, but consider whether to strengthen the theorem to the relaxed win condition (any consistent short key), which the same proof gives.}
